\section{本讲主线：并行训练是一张状态和通信的账本}

本讲不是并行术语清单，而是训练可行性分析：先从数据中心网络与 collectives 建立硬件边界，再依次切分 batch、训练状态、层、矩阵、序列、专家和上下文，最后把这些维度组合成真实模型配置。读每张图时都要追踪“每张卡长期保存什么、临时聚合什么、每步通信多少”。

\slidepair{slides-images/slide-002.jpg}{slides-images/slide-003.jpg}{课程目标与三部分组织：网络基础、并行 primitives、真实大模型组合策略。}{2--3}

Slide 2 的三个学习目标分别对应 systems complexity、parallelization paradigms 与 large-scale run configuration；Slide 3 把它们排成依赖顺序。若不了解网络拓扑，就无法判断 collective 成本；若不了解单一 primitive，就无法读懂最后的 3D/4D configuration。课程结构本身就是一条从硬件到模型 recipe 的推理链。

\begin{knowledgebox}{读图：Slide 1 为什么重要}
标题页看似只给题目，但它标记了课程从 Lecture 7 的通信 primitives 进入系统化并行训练。Lecture 7 让我们认识 rank、world size、collectives（集合通信，所有 rank 共同执行的多 GPU 通信原语，如 all-reduce、all-gather、reduce-scatter）、NCCL；Lecture 8 要回答的是，真实 LLM 训练中如何同时切分参数、梯度、优化器状态、激活、序列、专家和 batch。
\end{knowledgebox}


\begin{knowledgebox}{讲义补充：源 Slide 2 的三个学习目标}
第一，理解为什么大模型训练不是单卡训练的放大版，而是系统问题。第二，区分多种 parallelization paradigms：data、ZeRO/FSDP、pipeline、tensor、sequence、expert、context。第三，看懂大规模训练 run 的常见布局，也就是 DP/TP/PP/EP/CP 等数字如何组合成一张集群切分图。
\end{knowledgebox}


\begin{importantbox}{本讲的核心判断}
并行训练不是选择一个“最强技巧”。它是把训练状态放到不同设备上，并用 collective communication 在正确时刻恢复数学语义。读任何并行方案时，都要问四个问题：切了什么，复制了什么，通信什么，通信走哪条硬件路径。
\end{importantbox}

\begin{knowledgebox}{术语消化：本讲会反复出现的缩写}
\begin{tabular}{L{0.15\textwidth}L{0.33\textwidth}L{0.38\textwidth}}
\toprule
术语 & 定义 & 本讲关系 \\
\midrule
collectives & 集合通信，所有 rank 共同执行的多 GPU 通信原语，如 all-reduce、all-gather、reduce-scatter、broadcast。 & ZeRO、FSDP、TP、EP、CP 都能拆成这些 collective operations。 \\
sharding & 分片，把参数、梯度、optimizer state、activation 或 token 序列切到多张 GPU 上。 & 分片减少单卡显存，但通常增加 all-gather、reduce-scatter 或 all-to-all。 \\
ZeRO & Zero Redundancy Optimizer，按 stage 分片 optimizer state、gradients 和 parameters。 & 解决 data parallel 复制状态的显存浪费。 \\
FSDP & Fully Sharded Data Parallel，PyTorch 风格的 ZeRO-3 实现。 & 按模块临时 all-gather 参数，反向后 reduce-scatter 梯度。 \\
HBM & High Bandwidth Memory，高带宽显存，是 GPU 上的 DRAM。 & 训练状态最终要落在每张 GPU 的 HBM 里，显存容量决定是否需要分片。 \\
\bottomrule
\end{tabular}
\end{knowledgebox}

